\documentclass[11pt]{article}
\usepackage{geometry,hyperref}\geometry{margin=0.8in}
\title{Feedback and concurrency: a short research trail}
\date{15 November 2026}\begin{document}\maketitle
\noindent\textbf{Target.} Extend the existing single context/content construction with general feedback and concurrent interaction, while retaining native directed comparisons and pure endpoint changes. State threading, examples and a list of related theories do not by themselves meet that target. This is the current search trail, not a claim about the chronology of the original investigation.

\paragraph{1. Preserve the actual unifier.}
\textbf{Requirement:} keep the original operational and native-cell recoveries.\newline
\textbf{Sources:} the project's unified manuscript; Lawler, \url{https://arxiv.org/abs/1502.08017}; Capucci--Myers, \url{https://arxiv.org/abs/2410.21889}.\newline
\textbf{Contribution:} the existing base is a category-valued relative pseudomonad with a specified pure inclusion, not an ordinary Set monad. The original construction already covers its stated Kleisli/coKleisli, Para, span/content-dual and further algebra specializations.\newline
\textbf{Narrowing:} extend that retained presentation. Do not replace it with a hom-set quotient, or require all companions/conjoints merely because a more restrictive equipment formalism uses them.

\paragraph{2. Derive structure rather than assume the mixed law.}
\textbf{Requirement:} independently meaningful context/content hypotheses.\newline
\textbf{Sources:} existing O1; von Glehn, \url{https://www.tac.mta.ca/tac/volumes/33/36/33-36abs.html}; Capucci--Myers structured-contextad results.\newline
\textbf{Contribution:} display and structured-ambient methods organize universal-property proofs. The structured-paradise/realization theorems were already v1 Theorems 5.10/5.16, not discoveries in v2.\newline
\textbf{Narrowing/result:} the reviewed display-exact criterion derives the crossing and its coherence. The subsequent reviewed strict structured version derives a structured mixed realization from primitive compatibility, rather than supplying a structured crossing as an assumption. Neither establishes arbitrary physical interference.

\paragraph{3. Put feedback at the correct categorical level.}
\textbf{Requirement:} an operation on mixed process hom categories, not merely on native comparison arrows.\newline
\textbf{Source:} Hu--Vicary, \url{https://arxiv.org/pdf/2112.14051}.\newline
\textbf{Contribution:} traced pseudomonoids externalize trace using profunctor structure.\newline
\textbf{Narrowing:} ordinary Prof is a template; its Cat-valued relative adaptation is not proved by citation. Likewise a traced-category algebra on the category of native arrows alone does not give horizontal process feedback.

\paragraph{4. Retain state comparisons instead of quotienting them.}
\textbf{Requirement:} a feedback construction compatible with directed native cells.\newline
\textbf{Source:} Di Lavore et al., structured state feedback, \url{https://www.marioroman.xyz/papers/spangraph-open-transition.pdf}, Section 5.\newline
\textbf{Contribution:} its free state construction uses a coend and sliding along state maps, without a yanking axiom. The ordinary coend identifies those representatives.\newline
\textbf{Narrowing/current work:} use directed balance generators in a category-valued construction instead. The hom-category universal property and faithful native inclusion have been checked; the latter uses a state-intertwiner detector. This is not the category of all intertwiners, and inclusion need not be full. Sequential weak coherence is being checked; the complete free feedback doctrine remains unfinished. Stateful closure is not an instantaneous fixed-point solution.

\paragraph{5. Do not force effects to commute with time.}
\textbf{Requirement:} recursive/infinite behavior without invented infinite laws.\newline
\textbf{Sources:} Levy--Goncharov, \url{https://doi.org/10.4230/LIPIcs.CALCO.2019.13}; M\o gelberg--Zwart, \url{https://lmcs.episciences.org/16664/pdf}.\newline
\textbf{Contribution:} coinductive resumptions interleave effects with recursion; delay/effect distributive laws have substantive restrictions and obstructions.\newline
\textbf{Narrowing:} ordinary resumption existence does not supply a native relative crossing. The proved continuous final-sequence construction remains conditional on actual categorical continuity; no general trace or infinite effect law is inferred.

\paragraph{6. Contract shared grades coherently.}
\textbf{Requirement:} link several retained input/output grades without order-dependent identification.\newline
\textbf{Source:} Hirata, \emph{Notes on Lax Ends}, \url{https://arxiv.org/pdf/2210.01522}, Propositions 2.4 and 3.10.\newline
\textbf{Contribution:} lax coend Fubini exists, but requires cross-variable compatibility, not just separate lax naturality.\newline
\textbf{Narrowing:} use that established machinery for directed shared-grade contraction, keeping pure-middle invertible balance distinct. The current generator proof includes interchange and selected coordinate inversion. A walking-arrow test shows why replacing every pure coend by a directed one breaks the ordinary representable unit: its result is $0\to1$, not a terminal category. Thus pure horizontal balance remains invertible, while added shared-grade balance may be directed. Network evaluation is still separate; a wiring presentation alone is not unrestricted-concurrency semantics.

\paragraph{7. Keep unrestricted concurrency as a separate obligation.}
Shared-interface wiring, interleaving, joint dependence and sequential noncommutation are different structures. A symmetric tensor may model independent juxtaposition while interaction enters through wiring; it does not automatically certify arbitrary shared-state exchange. Premonoidal and network/equipment approaches are being compared, not imposed as a replacement. The outstanding task is a general composition/feedback structure admitting the intended interacting processes while deriving its compatibility with the existing retained context/content presentation.
\paragraph{8. Parallel composition without forcing commutativity.}
\textbf{Requirement:} genuine parallel assembly with directed, not necessarily invertible, sequential interchange.\newline
\textbf{Source:} Paquet--Saville, \emph{Effectful semantics in bicategories}, Definition 6.1 and Proposition 6.4, \url{https://philipsaville.co.uk/concurrent-pseudomonads.pdf}.\newline
\textbf{Contribution:} concurrent pseudomonads relax the multiplicative comparison of monoidal pseudomonads while retaining specified unit-whiskered invertibilities; their Kleisli parallel operation is a lax normal two-variable functor.\newline
\textbf{Narrowing:} investigate parallel structures on the independent context and content sides, together with a compatible crossing, that derive a directed interchanger on their relative composite. Our symmetric-monoidal state completion alone does not provide this structure. Compatibility with feedback and shared-interface admission remains to be proved; this source is not a ready-made relative unifier.

\paragraph{9. When directed balance recovers all intertwiners.}
\textbf{Requirement:} retain native state comparisons without quotienting them, and understand when the directed presentation has an explicit normal form.\newline
\textbf{Current derivation:} if the input state maps admit right adjoints in the process bicategory, adjunction units and directed slides suggest normalization to conjugation arrows $r_h f l_h^*\Rightarrow g$, equivalently state-intertwiner cells $r_hf\Rightarrow gl_h$.\newline
\textbf{Narrowing:} the stated normalization has passed independent written-proof review; no novelty claim is made. Only the selected state maps are assumed adjointable. It must not strengthen the original pure interface globally. Functional spans supply a recovery test, but a graph model does not replace the general proof or settle feedback/concurrency compatibility.
\end{document}
